\section{Local separation survives passage to the marked arithmetic field}
\label{sec:marked65}
The marked rank-three argument compares the Gaussian field $K$ with a
larger abelian CM field $L$. A pair of distinct characters on a local
group can become equal after restriction to a smaller group. We show
that the particular $L$ in the host construction preserves enough of
the original residual characters to exclude every cyclotomic link.
This supplies a local hypothesis of descent even at weight $e=1$.

% BEGIN M65A1
\paragraph{M65A1. A separating element in each decomposition group of the chosen field.}
Fix the source host data, including $\ell,N,e,\psi$, a rational
algebraic closure $\Omega$, complex conjugation $\sigma$ and a
prolongation $B$ at $\ell$. Write $G=G_{\mathbb Q}$,
$\eta=\eta_4$, $\varepsilon$ for the full $\ell$-adic cyclotomic
character and $\Psi=\varepsilon^e\psi_{\rm Gal}$ with coefficients
in $O=c.O$. Write $\alpha=\overline\Psi$. Under the cohomological
hypothesis of AFE16, make the choices of that proof through LFC22,
including its field $L$, link, frame $\Phi$ and mark $g_1$. For
\emph{every} prolongation $P$ above $\ell$ there exists
\[
 s_P\in D_P\cap G_L,\qquad \varepsilon(s_P)=1,\qquad
                      \alpha(s_P)^2\ne1 .
\]
Here $D_P\subset G$ is the rational decomposition group. In particular
the element is available on the relative local group at $P$ over $L$.

Indeed, at $P$ take the corrected Frobenius of AFE3 and its square $x$
as in AFE10 and HSR3. The correction multiplies a Frobenius lift by
inertia to cancel its \emph{full} cyclotomic value. At $\ell$, this is
possible because the cyclotomic extension is totally ramified. The
dyadic conductor test in AFE9--10 gives
$\alpha(x)^2\ne1$, and $\eta(x)=1$. The chosen field before adjoining
the imaginary quadratic factor is cut out by the subgroup
\[
 U=H_0\cap H_\theta(w)\cap H_b\cap H_m,\qquad G_L=U\cap\ker\eta .
\]
Since $\ell\mid M=\ell N$, MFC1 gives
$D_P\cap\ker\eta=D_P\cap H_0$. Thus $x\in H_0$. The other three
subgroups are normal of indices $\ell,\ell^b,\ell^m$, respectively.
Set $d=\ell^{1+b+m}$ and $s_P=x^d$. Each of the three finite
quotients kills this power. It remains in $H_0\cap\ker\eta\cap D_P$
and its full cyclotomic value is still one. In characteristic $\ell$,
$(\alpha(x)^2)^d-1=(\alpha(x)^2-1)^d\ne0$.
The field and its original mark have already been fixed. The local
witnesses need not be the mark and depend on no ordinary point.
% END M65A1

% BEGIN M65A2
\paragraph{M65A2. Ordinary flags and abelian inertia over the Gaussian field.}
Fix any embedding of the source $K$ into $\Omega$. Its image is the
fixed field of $\ker\eta$. Let $A$ be any one of the finite coefficient
DVRs of $\mathrm{AutEnd3Alt}$, let $r$ be an absolute
$\mathrm{OrdPt}$ for the chosen data, and let $\rho$ have exactly its
matrices on $G_K$ over $E=\overline{\mathbb Q}_\ell$. On $G_K$ these
matrices multiply ordinarily, since their sign is positive.

At a place away from $\ell$, put $\chi=\Psi|_{G_K}$ with values in
$E^\times$. It is continuous. Its inertia image is finite, since
$\varepsilon$ is trivial there and $\psi_{\rm Gal}$ has finite image.
If an inertia element has $\chi=1$, the away condition of the
absolute ordinary point gives $r=1$ there. Thus the restriction of
$\rho$ to inertia factors through a finite cyclic subgroup of
$E^\times$. The matrix of its generator diagonalizes in characteristic
zero, with eigenvalues powers of the generator of that subgroup.
There are consequently integers $j_0,j_1,j_2$ such that on this
entire inertia group the characteristic polynomial at $h$ is
$\prod_i(X-\chi(h)^{j_i})$. The three comparison characters extend
continuously to the decomposition group by the same powers of $\chi$.
This is precisely $\mathrm{IneAb}$, also when the inertia image is
trivial.

At $P\mid\ell$, write $P=tB$ with $t\in G$. At $B$ the defining
ordinary basis puts $r$ in upper triangular form on the sign-positive
part of the decomposition group. For $h\in D_P\cap G_K$ one has
$t^{-1}ht\in D_B\cap\ker\eta$. For positive $t$, conjugating the
original flag by the matrix of $r(t)$ gives a flag for $\rho$ at
$P$. For negative $t$ the signed conjugation inserts
$X\mapsto J (X^{-1})^{\mathsf T}J^{-1}$, with the source
antidiagonal $J$. This automorphism preserves upper triangularity:
inverse transpose changes upper to lower and reversal changes lower
back to upper. Apply it also to the change-of-basis matrix before
conjugating by the matrix of $r(t)$. In both cases the three
diagonal entries are characters by triangular multiplication.
The three inertia values are $\Psi,1,\Psi^{-1}$ in order, since
inversion and reversal together fix this ordered triple. This proves
$\mathrm{Flagged}$ for every prolongation over every place above
$\ell$. Restricting the same flags supplies flags on $G_L$, since
$G_L\subseteq\ker\eta$.
% END M65A2

% BEGIN M65A3
\paragraph{M65A3. Unlinked flags over the upper field, including weight one.}
For the field of M65A1 every place of $L$ above $\ell$ satisfies the
original $\mathrm{Unl1}(\rho|_{G_L},w)$. Fix a prolongation $P$ at
$w$. At $s_P$ the residual condition on $r$ identifies its reduced
matrix with the diagonal base matrix. Its characteristic roots are
\[
                   \alpha(s_P),\quad 1,\quad\alpha(s_P)^{-1}.
\]
They are pairwise distinct in the residue field. The embedding of $O$
in $A$ induces an injection on residue fields: $A$ is finite integral
over $O$ and its maximal ideal contracts to the maximal ideal of $O$.
The integral characteristic polynomial at $s_P$ thus has three
distinct roots over $E$ as well. Indeed a monic polynomial with
separable reduction has unit discriminant. The values at $s_P$ of
the three characters $\kappa_i$ of the flag from M65A2 are exactly
these characteristic-zero roots. For $i\ne j$ the identity
$\kappa_i=\kappa_j\varepsilon$ would make two of them equal,
since $\varepsilon(s_P)=1$. For $i=j$ this identity would force
$\varepsilon=1$ on the entire decomposition group. The local field
$L_w$ is finite over $\mathbb Q_\ell$ and cannot contain all
$\ell$-power roots of unity. This rules out the last possibility.
The argument applies at every $P$, as $\mathrm{Unl1}$ requires.

It follows that $\mathrm{Keeps}(\rho,L)$ holds: at every $w\mid\ell$
the unlinked condition over $L$ already holds, so it holds whenever
the lower flag is unlinked. Twisting all three characters by the
same continuous character preserves these ratios on both fields.
In particular this result applies to the inverse-cyclotomic twist.
The calculation uses the special field in M65A1. It establishes
the local preservation required in this host argument, without an
assumption about preservation on an arbitrary linked arithmetic frame.

There is a further consequence once the continuous twisted
representation over $L$ is supplied with $\mathrm{IsAutCsdA}$.
Its $\mathrm{WitA}$ first clause upgrades $\mathrm{Unl1}$ there to
$\mathrm{PotUnl}$, excluding the character identity even after
raising both sides to any common positive power. The same exclusion
then holds on $K$. For a fixed $P$, restrict the lower ordinary flag
to the upper decomposition group. Its one-dimensional constituents,
as a multiset, agree with those of the flag witnessing
$\mathrm{PotUnl}$ there by Jordan--H\"older. A power identity on the
lower flag would restrict to the same power identity for two of the
upper constituents. This contradicts the upper exclusion. Every
prolongation above a given place on $K$ restricts to one of the
covered places on $L$, proving the quantifier over all prolongations.
% END M65A3

% BEGIN M65A4
\paragraph{M65A4. The marked theorem from two upper-field assertions and polarized descent.}
For each original host take precisely the cohomological construction
of M65A1, with its own link $k$. Suppose its frame has the following
two global arithmetic conclusions at the fixed $\iota$:
\[
 \mathrm{FinPt}(O,\Phi,3),\qquad
 \mathrm{AutPtAlt}(k,3,\iota,\varepsilon_L^{-1}).
\]
Assume also the rank-three instance of the source polarized solvable
descent implication $\mathrm{DownCsdKAlt}$ for the resulting
extension $L/K$. These inputs imply $\mathrm{ThmS6}_s$ when supplied
for all original hosts. We verify the marked comparison.

For any embedding of $K$ into $\Omega$, $G_L\subseteq\ker\eta=G_K$
places its image inside the already chosen $L$ by the fixed-field
theorem. The induced embedding $K\to L$ gives the compatible scalar
tower in the source. The extension is finite abelian since
$L/\mathbb Q$ is finite abelian. AFE17 and AFE20 give the marking and
$\mathrm{FinEnd}$ at $g_1$ from the first global conclusion. In
particular the finite-index subgroup is chosen before every residual
DVR point.

Now take an eligible absolute ordinary $r$ and $\rho$ as in M65A2,
with the open-subgroup irreducibility required in
$\mathrm{AutEnd3Alt}$. Since $G_L$ is open in $G$ and contained in
$G_K$, this hypothesis gives irreducibility of $\rho|_{G_L}$.
AFE19 constructs a normalized rank-three representation over $A$
on the same frame. It is a lift of the required residual point.
At each specified inertia element its three diagonal characters are
the coefficient images of $\Psi^{1},1,\Psi^{-1}$. On $G_L$ its
matrices after coefficient inclusion equal those of $\rho|_{G_L}$
through the link map $k.p_0$. These are the lift, weights and
literal matrix identification in the second global conclusion.
It supplies $\mathrm{IsAutCsdA}$ at the original $\iota$ for
$\tau|_{G_L}$, where $\tau=\rho\varepsilon_K^{-1}$.

We check descent on $\tau$ itself. Entrywise adic continuity of $r$
implies $\ell$-adic continuity after embedding $A$ in $E$, because
a uniformizer of $A$ has absolute value less than one. The
cyclotomic factor is continuous. Twisting preserves invariant
subspaces and multiplies all the local characters by the same
factor, so irreducibility on $G_L$, abelian inertia, the flags and
the local preservation follow from M65A2--3. For polarization,
normalization at $\sigma$ gives matrix $1$ and sign $-1$.
The signed group law consequently gives, whenever
$s'=\sigma s\sigma^{-1}$ for $s,s'\in G_K$,
\[
                   \rho(s)^{\mathsf T}J\rho(s')=J .
\]
Since $\varepsilon(s')=\varepsilon(s)$, multiplying both
matrices by their inverse-cyclotomic factors gives
$\tau(s)^{\mathsf T}J\tau(s')=\varepsilon(s)^{-2}J$.
This is exactly the required multiplier $\varepsilon^{1-3}$
with the specified conjugation. The descent implication now
supplies the full $\mathrm{IsAutCsdA}(\iota,\tau)$ over $K$,
including its window for every attached cusp. This proves
$\mathrm{AutEnd3Alt}$ for the same $g_1$ as the finite end.
VMA6 gives the stated marked theorem.

The two upper-field assertions and the indicated polarized descent
must still be supplied to deduce the whole obligation from its
original hypotheses. The local comparison above already covers
all weights, both Gaussian embeddings and every eligible ordinary point
on the field chosen by the host construction.
% END M65A4
